\chapter{Cross-tome closure audit: связь Томов I--III}

Полный контекст первых двух томов требует отделить локальные успехи
версии III от исходной программы observed-world unification.

\section{Окончательная проверка factor-two dispute}

Для одной координаты достаточно записать
\begin{equation}
 L=A\dot q^2-\frac12Hq^2.
\end{equation}
Euler--Lagrange equation равна
\begin{equation}
 2A\ddot q+Hq=0,
\end{equation}
поэтому
\begin{equation}
 m^2=\frac{H}{2A}.
\end{equation}
В нашем случае
\begin{equation}
 A=\frac{\kappa}{2}G,
\end{equation}
и, следовательно,
\begin{equation}
 m^2=(\kappa G)^{-1}H.
\end{equation}
При \(\kappa=2\) каноническая matrix равна \(2G\), а не \(G\).
Контекст Тома II не меняет это variational identity.

\begin{theorem}[Final kinetic-normalization verdict]
Предложенный пересчёт \(40\to184\) и \(67\to211\) отклоняется.
Корректными остаются
\begin{equation}
 m_s^2=4\chi^2,
 \qquad
 N_{\mathrm{fin}}=40,
 \qquad
 N_{\mathrm{fin+gauge}}=67.
\end{equation}
Разногласие вызвано использованием coefficient \(A\) как mass metric без
обязательного множителя \(2\) в Euler--Lagrange equation.
\end{theorem}

\section{Статус \(S_{\mathrm{vac}}\) Тома II}

Том II окончательно классифицирует
\begin{equation}
 S_{\mathrm{vac}}
 =
 S_{\mathrm{geo}}
 -\frac1{24S_{\mathrm{geo}}}
 -\frac1{\pi^4S_{\mathrm{geo}}^2}
\end{equation}
как strong structural compression. Periodic Maxwell--ghost branch
\(1/24\) сохраняет условный determinant status, но C6 closure для
\(\pi^{-4}\)-члена закрыта отрицательно: full same-scheme compensation
не найдена, а Casimir kernel не совпадает с Euclidean winding
log-determinant.

Версия III строит другой объект:
\begin{equation}
 \Gamma_{\mathrm{hidden}}[\Phi,A,\chi].
\end{equation}
Она не выводит \(S_{\mathrm{vac}}\) как stationary value, coupling,
counterterm или determinant residue.

\begin{nogocriterion}[Missing \(S_{\mathrm{vac}}\)--CW bridge]
Численное соседство геометрических coefficients Тома II и
Coleman--Weinberg scale Тома III не считается unified mechanism.
Необходим один functional \(\Gamma_{\mathrm{unif}}\), две контролируемые
редукции которого дают одновременно \(S_{\mathrm{vac}}\) и
\(\Gamma_{\mathrm{hidden}}\).
\end{nogocriterion}

\section{EW/QCD verdict}

Том II закрывает минимальные EW/QCD repair branches отрицательно:
\begin{enumerate}
 \item ordinary split running имеет неправильное направление;
 \item logarithmic threshold cone не содержит физического решения;
 \item two-loop stress test оставляет
 \(\alpha_s(M_Z)=0.080177\), то есть примерно \(-32\%\) относительно
 используемого контрольного значения;
 \item nonlogarithmic matching mechanism не выведен.
\end{enumerate}

Hidden relative-\(U(1)\) action Тома III не содержит
\(SU(2)\), \(SU(3)\), Standard Model hypercharge representation или
weak-current vertices. Поэтому positive \(B\), scale-setting и hidden
vacuum не исправляют EW/QCD no-go.

\begin{theorem}[Scale gate is not physical closure]
Даже полное замыкание hidden Coleman--Weinberg scale было бы только
необходимым условием внутренней согласованности версии III. Оно не
является достаточным условием воспроизведения observed gauge sector.
\end{theorem}

\section{Два появления числа \(23\)}

В Томе II
\begin{equation}
 23+\pi^{-1}
\end{equation}
является условной collective neutrino norm. Rank-one selector в полном
generation--spinor module не выведен существующими симметриями; defect
route также требует дополнительных operator data.

В раннем finite supertrace Тома III число \(23\) появлялось как
\begin{equation}
 184=8\cdot23
\end{equation}
в иной kinetic convention. После heat-kernel normalization актуальный
finite numerator равен \(40\). Ни прежнее, ни текущее arithmetic не
создают neutrino rank theorem.

\begin{proposition}[No structural \(23\)-bridge]
Между neutrino collective rank claim и finite scalar supertrace не
построено representation identity, index theorem или common projector.
Оба контекста должны оставаться логически раздельными.
\end{proposition}

\section{Orbit half-trace}

Правило
\begin{equation}
 \Tr_{\mathrm{orb}}=\frac12\Tr_{H_8}
\end{equation}
устраняет формальное KO6 particle--antiparticle doubling и согласованно
используется во всех finite calculations версии III. Однако ни Том I, ни
Том II, ни текущая functional measure не выводят его как обязательный
bosonic spectral trace.

\begin{nogocriterion}[Cross-tome half-trace gap]
Hidden parent action считается полностью выведенным только если
orbit half-trace возникает из Pfaffian, quotient measure, groupoid trace
или другой единой measure construction. До этого absolute multiplicities
и gauge-normalization relation имеют условный measure status.
\end{nogocriterion}

\section{Роль Тома I}

Том I задаёт falsification methodology, но не операторную модель.
Текущее состояние удовлетворяет его критерию незавершённости:
\begin{equation}
 \text{единый observed-world parent functional не построен}.
\end{equation}
Это не является логическим опровержением онтологической гипотезы, но
запрещает называть три тома завершённой unified physical theory.

\section{Сводный статус}

\begin{center}
\begin{tabular}{p{0.29\textwidth}p{0.20\textwidth}p{0.33\textwidth}}
\toprule
Узел & Статус & Точный смысл \\
\midrule
геометрическое ядро \(K\) & условно сохранено & не доказана неизбежность carrier \\
\(S_{\mathrm{vac}}\) & structural compression & C6 determinant closure negative \\
neutrino \(23+\pi^{-1}\) & conditional & selector/readout open \\
EW/QCD & negative minimal branch & correct-sign and threshold gates fail \\
Version III parent block & mathematical pass & hidden \(U(1)\) EFT \\
absolute scale & one train datum & spectral moment orbit open \\
orbit half-trace & conditional & measure origin open \\
observed-world unification & not achieved & common functional absent \\
\bottomrule
\end{tabular}
\end{center}

\begin{theorem}[Cross-tome verdict]
Версия III является реальным прогрессом относительно Тома II: построен
конкретный anomaly-free hidden parent action с выведенными relative
normalizations. Но он не закрывает исходную S2T-программу observed
microworld, поскольку не связан единым functional с \(S_{\mathrm{vac}}\),
neutrino operator и EW/QCD sector.
\end{theorem}

\section{Разделение дальнейших программ}

Дальнейшая работа должна быть разделена:
\begin{enumerate}
 \item \textbf{III.H --- hidden completion:} вывести orbit measure и
 portal, после чего сформулировать реальные hidden-sector observables;
 \item \textbf{IV.SM:} построить observed-sector finite
 algebra/representation с
 \(SU(3)\times SU(2)\times U(1)\), neutrino graph и common bridge к
 electromagnetic functional.
\end{enumerate}
Запрещено добавлять Standard Model sectors в hidden action по одному и
считать это продолжением того же минимального parent block.

\begin{protocol}[Следующий приоритет]
В рамках Тома III минимальный незакрытый внутренний узел --- происхождение
orbit half-trace из functional measure. Portal имеет смысл строить только
после measure closure. Observed-sector reconstruction относится уже к
отдельной версии и не должно смешиваться с текущим hidden scorecard.
\end{protocol}

Машинный cross-tome ledger сохранён в
\path{s2t_v3_cross_tome_closure_results.json}.